\section{Experiments}
\label{sec:experiments}

\subsection{Datasets and reproduction}
\label{sec:datasets}
We train and evaluate in-domain on \textbf{STARSS23} (DCASE~2024 Task~3
development set), which provides both MIC and FOA formats with $13$ sound
classes. For external validity we evaluate zero-shot on \textbf{STARSS22}
(DCASE~2022 Task~3 development-test, $54$ clips; distance-blind matching, as
STARSS22 lacks distance labels) and stress-test on the synthetic
\textbf{TAU-NIGENS-SSE-2021} corpus, remapping the seven classes that overlap
with the STARSS taxonomy. Before ablating we confirm pipeline parity: the FOA
reproduction (five seeds) gives $F_{20\degsym}=13.21\pm0.63\%$ against the
published $13.10\%$, with DOAE$_\mathrm{CD}=41.5\pm5.5\degsym$ and
RDE$=0.27\pm0.01$, inside challenge variance. Subsequent differences are thus
attributable to the toggled factors, not a mis-implemented baseline.
\label{sec:repro}

\subsection{Main result: the effect is graded by inductive bias}
\label{sec:main}
Table~\ref{tab:dissociation} reports the headline contrast,
$\Delta\text{DOAE}_\mathrm{CD}=\text{full}-\text{no\_geom}$, seed-matched within
each cell ($n=5$) and grouped by temporal backbone; Fig.~\ref{fig:dissociation}
visualizes it. The effect on direction follows a \textbf{consistent ordering set
by the backbone's inductive bias}: \emph{both recurrent cells shift
non-positively} ($\Delta\le0$; FOA a clear benefit, MIC negligible at
$-0.9\degsym$, $\dz=-0.07$), \emph{both Transformer cells shift positively}
($\Delta\ge0$), and the convolution--attention Conformer sits between them
(helpful on MIC, neutral on FOA). Averaged over modality the per-backbone mean $\Delta$ moves monotonically
from $-2.9\degsym$ (CRNN) through $-0.5\degsym$ (Conformer) to $+4.1\degsym$
(Transformer) as temporal locality is removed. The two extreme effects sit on
opposite ends of this ordering: the prior improves localization by $5.0\degsym$
in FOA+CRNN ($\dz=-1.03$) and degrades it by $5.7\degsym$ in MIC+Transformer
($\dz=+0.71$).

At $n=5$ the individual per-cell contrasts are trends rather than significant in
isolation (e.g.\ FOA+CRNN $p=0.082$, MIC+Transformer $p=0.186$). \emph{These
single-cell $p$-values are not our primary evidence}; the claim is about the
\emph{ordering} across the grid, which we test jointly with a factorial
interaction (Section~\ref{sec:anova}) and corroborate on an independent dataset
where the recurrent benefit reaches full significance
(Section~\ref{sec:cross}).

One honest caveat about the middle of the ordering: the Conformer is the weakest
detector in our grid (location-aware $F_{20\degsym}$ of $4$--$7\%$ and the worst
SELD scores, vs.\ $9$--$13\%$ for the extremes), so its near-zero geometry effect
may partly reflect a weaker operating point rather than intermediate locality. We
therefore treat it as \emph{corroborating} the ordering, not proving it: the
load-bearing comparison is the recurrent-versus-Transformer contrast between two
healthy detectors, with the trend test (Section~\ref{sec:anova}) confirming only
that the hybrid does not violate the ordering.

\begin{table}[t]
\centering
\caption{Main result, grouped by backbone: paired geometry-prior contrast on
class-dependent localization error,
$\Delta\text{DOAE}_\mathrm{CD}=\text{full}-\text{no\_geom}$ (degrees; negative
$=$ prior helps), $n=5$ each. The per-backbone mean sign ordering CRNN $\le$
Conformer $\le$ Transformer holds after averaging over modality; the recurrent
and Transformer extremes also order as predicted within each modality.
Source: \texttt{path\_c\_2x2}.}
\label{tab:dissociation}
\setlength{\tabcolsep}{4pt}
\begin{tabular}{llccc}
\toprule
Backbone & Modality & $\Delta$DOAE & $t$ ($p$) & \dz \\
\midrule
\multicolumn{5}{l}{\emph{CRNN+MHSA (recurrent): prior helps}}\\
CRNN        & MIC & $-0.88$ & $-0.15$ ($0.887$) & $-0.07$ \\
\textbf{CRNN} & \textbf{FOA} & $\mathbf{-5.02}$ & $-2.31$ ($0.082$) & $\mathbf{-1.03}$ \\
\midrule
\multicolumn{5}{l}{\emph{Conformer (conv--attention hybrid): intermediate}}\\
Conformer   & MIC & $-1.83$ & $-0.35$ ($0.745$) & $-0.16$ \\
Conformer   & FOA & $+0.88$ & $+0.30$ ($0.778$) & $+0.14$ \\
\midrule
\multicolumn{5}{l}{\emph{Transformer-only: prior hurts}}\\
\textbf{Transformer} & \textbf{MIC} & $\mathbf{+5.70}$ & $+1.60$ ($0.186$) & $\mathbf{+0.71}$ \\
Transformer & FOA & $+2.57$ & $+0.86$ ($0.437$) & $+0.39$ \\
\bottomrule
\end{tabular}
\end{table}

\begin{figure}[t]
\centering
\includegraphics[width=\columnwidth]{path_c_2x2_dissociation.png}
\caption{The geometry-prior effect on DOAE$_\mathrm{CD}$ across the
modality$\times$backbone grid. The sign of the effect tracks temporal inductive
bias: it helps (negative) under recurrence, is neutral for the Conformer, and
hurts (positive) under the Transformer, the extremes being FOA+CRNN and
MIC+Transformer.}
\label{fig:dissociation}
\end{figure}

\subsection{Primary test: architecture$\times$prior interaction}
\label{sec:anova}
Our primary hypothesis test asks whether the sign ordering of
Table~\ref{tab:dissociation} is a genuine interaction rather than a collection of
independent per-cell coincidences. We report it three complementary ways on the
per-seed DOAE$_\mathrm{CD}$, all pointing the same direction
(Source: \texttt{path\_c\_trend\_mixed}).

\emph{(i) Factorial interaction, seed-matched.} A Type-II factorial ANOVA
restricted to the recurrent and Transformer extremes
(modality$\times$architecture$\times$prior, $n_{\text{obs}}=40$) gives a
significant \textbf{architecture$\times$prior interaction} ($F=4.62$, $p=0.039$):
the prior's effect on direction depends on whether the backbone is recurrent or
pure-attention. Because the same five seeds are matched across cells, we re-express
this interaction as a within-seed paired \emph{second difference},
$D=\Delta_{\text{CRNN}}-\Delta_{\text{Xfm}}$ per matched (modality, seed), where
$\Delta=\text{full}-\text{no\_geom}$; $D$ is reliably negative
($\bar D=-7.1\degsym$, $t=-2.59$, $p=0.029$, $\dz=-0.82$), i.e.\ the prior helps the
recurrent backbone about $7\degsym$ more than it helps the Transformer. A linear
mixed model with a per-seed random intercept estimates that variance at the zero
boundary (ICC$\approx0$ for DOAE$_\mathrm{CD}$), so the design carries no extra
between-seed structure and the fixed-effects estimate is the appropriate one.

\emph{(ii) Monotone trend across all three backbones.} Rather than discard the
intermediate Conformer, we test the \emph{ordered} hypothesis that the geometry
effect $\Delta$ increases monotonically along the locality axis
CRNN$<$Conformer$<$Transformer. A Jonckheere--Terpstra trend test on the per-seed
$\Delta$ (pooled over modality) is a one-sided trend at $z=+1.60$, $p=0.055$, with a
parametric slope of $+3.5\degsym$ per locality step ($p=0.073$); the trend is
clearest in Ambisonics (one-sided $p=0.051$) and is the weakest in the microphone
array, where the Conformer is itself mildly helpful. The monotone ordering of the
per-backbone means ($-2.9\degsym\!\to\!-0.5\degsym\!\to\!+4.1\degsym$) is therefore
a genuine---if, at $n=5$, marginal---trend, not an artifact of selecting the two
extremes.

\emph{(iii) Unordered omnibus.} When the Conformer is added as a third level in the
\emph{unordered} factorial ($n_{\text{obs}}=60$), the three-level
architecture$\times$prior term is not significant ($F=1.94$, $p=0.15$), exactly as
expected when a near-zero middle category is inserted between two opposing extremes:
an unordered omnibus spends power on a quadratic contrast the data do not exhibit,
whereas the ordered trend test in (ii) retains it. The prior main effect is null
throughout ($p>0.6$), confirming the effect is conditional rather than a universal
benefit or harm. (On the detection metric $F_{20\degsym}$ the design shows strong
modality $F=36.9$ and architecture $F=79.2$ main effects, consistent with FOA and
recurrence detecting events more readily overall; these concern detectability, not
the prior.)

\subsection{Anchor: the ordering replicates zero-shot on STARSS22}
\label{sec:cross}
The grid ordering is not a STARSS23 artifact. Evaluating the STARSS23-trained
checkpoints zero-shot on STARSS22 (Table~\ref{tab:cross} and
Fig.~\ref{fig:forest}, $n=5$) reproduces the same architecture ordering
out-of-domain: both recurrent cells help, the Conformer is helpful-to-neutral,
and both Transformer cells hurt. Fig.~\ref{fig:forest} makes the agreement
visible---in-domain and zero-shot estimates share the same sign in every cell,
and the effect marches monotonically from the ``helps'' to the ``hurts'' region
as locality is removed. The single
\emph{significant} cross-dataset effect is the recurrent benefit on Ambisonics,
FOA+CRNN $\Delta\text{DOAE}_\mathrm{CD}=-4.76\degsym$ ($\dz=-1.50$, $p=0.029$),
which matches its in-domain sign and magnitude ($-5.02\degsym$)---a benefit
reproducing on an independent recording campaign is strong evidence that it is a
property of the learned model, not of one test set. The MIC+Transformer harm
reproduces directionally ($+5.58\degsym$, $\dz=+0.77$) but, like in-domain, does
not reach significance on its own at $n=5$. The cross-dataset per-backbone means
trace the same monotonic path as in-domain (helpful $\to$ neutral $\to$ harmful).

\begin{table}[t]
\centering
\caption{Anchor. STARSS22 zero-shot, geometry-prior contrast on
DOAE$_\mathrm{CD}$ (degrees; negative $=$ helps), $n=5$. The in-domain ordering
holds out-of-domain; the FOA+CRNN benefit is significant.
Source: \texttt{path\_c\_cross\_starss22}.}
\label{tab:cross}
\setlength{\tabcolsep}{4pt}
\resizebox{\columnwidth}{!}{%
\begin{tabular}{llcc}
\toprule
Backbone & Modality & $\Delta$DOAE (cross) & in-domain \\
\midrule
CRNN        & MIC & $-4.77$ ($\dz=-0.42$) & $-0.88$ \\
\textbf{CRNN} & \textbf{FOA} & $\mathbf{-4.76}$ ($\dz=-1.50$, $p\!=\!.029$) & $-5.02$ \\
Conformer   & MIC & $-3.20$ ($\dz=-0.48$) & $-1.83$ \\
Conformer   & FOA & $-0.42$ ($\dz=-0.06$) & $+0.88$ \\
Transformer & MIC & $+5.58$ ($\dz=+0.77$) & $+5.70$ \\
Transformer & FOA & $+1.72$ ($\dz=+0.39$) & $+2.57$ \\
\bottomrule
\end{tabular}}
\end{table}

\begin{figure}[t]
\centering
\includegraphics[width=\columnwidth]{forest_gradient.png}
\caption{Geometry-prior effect on directional error
($\Delta\text{DOAE}_\mathrm{CD}=\text{full}-\text{no\_geom}$; negative $=$ helps)
per cell, ordered by temporal locality, with bootstrap $95\%$ CIs. Filled circles
are in-domain (STARSS23); open squares are zero-shot on STARSS22; ``$*$'' marks
$p<0.05$. The effect shifts monotonically from the helps region (recurrent) to
the hurts region (Transformer), and the in-domain and zero-shot estimates agree
in sign for every cell---the recurrent Ambisonic benefit being significant
out-of-domain.}
\label{fig:forest}
\end{figure}

\subsection{Mechanism probe: the effect is behavioural, not a linear feature shift}
\label{sec:probing}
A natural hypothesis is that the prior changes how linearly direction is encoded
in the shared post-convolutional representation. We test this by linearly probing
those features (frequency-pooled, ridge regression to $(\sin,\cos)$ of azimuth
and elevation on single-source frames, file-wise five-fold CV) and comparing
\texttt{full} against \texttt{no\_geom} (Table~\ref{tab:probe}). At $n=5$,
\textbf{no cell shows a significant change} in linearly decodable direction
($|\dz|\le0.53$, all $p>0.2$); the largest movement is a non-significant increase
in MIC+Conformer. We report this as an informative negative result: the prior's
behavioural effect on SELD output is \emph{not} mirrored by a measurable shift in
the linear direction content of the post-conv features, so the mechanism does not
reduce to ``the prior adds/removes decodable direction information.'' This null is
non-trivial. The geometry bias gates the input channels \emph{upstream} of the
shared convolutional front-end, and the \texttt{full} and \texttt{no\_geom}
variants train their own convolutional weights, so the probed post-convolutional
features could in principle diverge; that they do not differ in linearly decodable
direction means the prior leaves the linearly accessible spatial content of the
shared representation essentially unchanged. Its behavioural effect must therefore
reside in nonlinear or distributed feature changes the linear probe cannot
capture, or in how each temporal stack \emph{uses} that representation---the
latter consistent with the architecture-dependence we observe. (We caution that an earlier three-seed probe suggested a
Transformer-specific loss; it did not survive the move to five seeds, a useful
reminder of the fragility of small-$n$ representational claims.)

\begin{table}[t]
\centering
\caption{Mechanism probe. Linear-probe angular MAE (degrees; lower $=$ direction
more linearly decodable), $\Delta=\text{full}-\text{no\_geom}$, $n=5$. No cell
shows a significant geometry effect. Source: \texttt{path\_c\_probe}.}
\label{tab:probe}
\setlength{\tabcolsep}{4pt}
\begin{tabular}{llccc}
\toprule
Backbone & Modality & full & no\_geom & $\Delta$ \\
\midrule
CRNN        & MIC & $28.50$ & $28.37$ & $+0.12$ \\
CRNN        & FOA & $20.62$ & $20.89$ & $-0.27$ \\
Conformer   & MIC & $29.88$ & $28.35$ & $+1.53$ \\
Conformer   & FOA & $18.69$ & $19.28$ & $-0.59$ \\
Transformer & MIC & $26.11$ & $26.27$ & $-0.16$ \\
Transformer & FOA & $17.93$ & $17.89$ & $+0.05$ \\
\bottomrule
\end{tabular}
\end{table}

\subsection{Boundary conditions and supporting ablations}
\label{sec:synth}
\textbf{Distance follows a different rule.} Beyond direction, the prior's effect
on the third (distance) dimension does \emph{not} track the temporal architecture
(Table~\ref{tab:rde}): relative distance error improves by effect size in
\emph{all three} Ambisonic cells---FOA+CRNN $\Delta\text{RDE}=-0.024$
($\dz=-1.43$, $p=0.033$), FOA+Transformer $-0.020$ ($\dz=-1.00$), FOA+Conformer
$-0.010$ ($\dz=-0.35$)---and is null in all three microphone-array cells
($|\dz|\le0.26$, all $p>0.5$), with no architecture$\times$prior interaction
($F=0.20$, $p=0.82$; paired CRNN$-$Xfm second difference $p=0.52$). The
architecture that governs the \emph{direction} effect is thus inert on distance,
where the benefit is instead concentrated in Ambisonics. We are careful not to
overclaim the second half of this contrast: at $n=5$ the modality$\times$prior
interaction on distance is itself only a non-significant trend ($F=0.39$,
$p=0.54$; paired FOA$-$MIC second difference $p=0.54$), so the distance pattern is
an effect-size dissociation---architecture selectively modulates direction and not
distance---rather than a second statistically confirmed interaction.
Source: \texttt{path\_c\_dissociation\_test}.

\begin{table}[t]
\centering
\caption{Distance is modality-patterned, not architecture-conditional.
Geometry-prior contrast on relative distance error,
$\Delta\text{RDE}=\text{full}-\text{no\_geom}$ ($n=5$; negative $=$ helps). All
Ambisonic cells improve by effect size; all microphone-array cells are
null---orthogonal to the architecture-conditional direction effect of
Table~\ref{tab:dissociation}, though the modality$\times$prior interaction is
itself only a non-significant trend at $n=5$ (see text).
Source: \texttt{path\_c\_2x2}.}
\label{tab:rde}
\begin{tabular}{llcc}
\toprule
Backbone & Modality & $\Delta$RDE & \dz \\
\midrule
CRNN        & MIC & $-0.016$ & $-0.26$ \\
\textbf{CRNN} & \textbf{FOA} & $\mathbf{-0.024}$ ($p\!=\!.033$) & $\mathbf{-1.43}$ \\
Conformer   & MIC & $-0.008$ & $-0.17$ \\
Conformer   & FOA & $-0.010$ & $-0.35$ \\
Transformer & MIC & $-0.002$ & $-0.04$ \\
Transformer & FOA & $-0.020$ & $-1.00$ \\
\bottomrule
\end{tabular}
\end{table}

\textbf{Boundary condition (synthetic transfer).} On the synthetic TAU-NIGENS
corpus (seven mapped classes, zero-shot; CRNN and Transformer cells, $n=3$--$5$)
detection collapses to near-floor for every cell (kept-class
$F\!\approx\!1$--$2\%$, DOAE$_\mathrm{CD}\!\approx\!60\degsym$) under the heavy
synthetic domain shift, and at that floor the geometry-prior contrasts are
uniformly null. Notably the in-domain MIC+Transformer harm does \emph{not}
transfer: the synthetic $\Delta\text{DOAE}_\mathrm{CD}$ is $+1.1\degsym$
($\dz=+0.17$) versus $+5.7\degsym$ in-domain, and the MIC+CRNN and FOA contrasts
are likewise null ($|\dz|\le0.6$). We read this as a boundary condition: a domain
shift severe enough to erase detection also erases the subtle prior effect,
whereas the real-data architecture ordering (Sections~\ref{sec:main}--\ref{sec:cross})
does not depend on it.

\textbf{Supporting controls.} A per-class breakdown ($13$ classes, all cells)
shows detectable $F$-score concentrating in speech, music, and domestic-sound
classes, so the architecture ordering is a property of the localization
representation, not of a few classes' detectability. A vanilla
squeeze-and-excitation control in MIC+CRNN confirms that the \emph{geometric}
bias, not channel attention itself, drives the contrast.
